\documentclass[../../main.tex]{subfiles}
\graphicspath{{\subfix{../../image/}}}

\begin{document}

\section{数学期望}\label{section:ch04-数学期望}

\subsection*{一、平均值与加权平均值}

有甲、乙两个射手，他们的射击技术用下表表出：

\begin{center}
\begin{tabular}{c|ccc}
\multicolumn{4}{c}{甲射手}\\
击中环数 & 8 & 9 & 10\\
\hline
概\ 率 & $0.3$ & $0.1$ & $0.6$
\end{tabular}
\qquad
\begin{tabular}{c|ccc}
\multicolumn{4}{c}{乙射手}\\
击中环数 & 8 & 9 & 10\\
\hline
概\ 率 & $0.2$ & $0.5$ & $0.3$
\end{tabular}
\end{center}

试问哪一个射手技术较好？

这个问题的答案不是一眼看得出的。这说明分布列虽然完整地描述了随机变量，但是却不够``集中''地反映出它的变化情况。因此我们有必要找出一些量来更集中、更概括地描述随机变量，这些量多是某种平均值。

求平均值是大家都很熟悉的一种运算。例如，某公司有 $n$ 个职工，他们的工资分别为 $x_1,x_2,\cdots,x_n$，则这个公司的平均工资为
\begin{align}
\bar{x}=\frac{x_1+x_2+\cdots+x_n}{n}.\label{align:ch04-算术平均}
\end{align}

还有其他求平均值的方法。例如，一个小学生，他的考试成绩为：语文 $95$ 分，算术 $85$ 分，常识 $60$ 分。若依上面方法计算，则他的平均成绩为 $\bar{x}=80$。显然，这个数字不太能反映这个学生的真正成绩，因为它没有考虑到这三个科目的相对重要性。在这个年级中，每周语文有 $10$ 节课，数学有 $8$ 节，而常识只有 $2$ 节。在评价学生成绩时，这个因素不能不考虑，因此用下面方法来计算学生的平均成绩，似乎更合理些：
\begin{align*}
\bar{x}_w&=\frac{95\times 10+85\times 8+60\times 2}{10+8+2}\\
&=95\times\frac{10}{20}+85\times\frac{8}{20}+60\times\frac{2}{20}=87.5.
\end{align*}

这种平均称为\textbf{加权平均}。其一般定义如下：给定权 $w_i\geqslant 0$，$i=1,2,\cdots,n$，满足 $\sum_{i=1}^n w_i=1$，则
\begin{align}
\bar{x}_w=\sum_{i=1}^n w_i x_i\label{align:ch04-加权平均}
\end{align}
称为 $x_1,x_2,\cdots,x_n$ 关于权 $\{w_i,\,i=1,2,\cdots,n\}$ 的\textbf{加权平均值}。

显然，在某些情况下，加权平均更加合理。由于``权''的大小直接影响最后结果，因此``权''的选择是加权平均中最重要的问题。例如为测量泰山的高度可以安排几个测量队从不同的地点进行测量，由于各个测量队的技术水平不一，各地点的地理条件不同，因此最后的结果恐怕要采用某种加权平均值为好。

不言而喻，普通平均是加权平均的一种特例，这时所有的权相等。

平均值按其大小总在原始数据的当中，因此它反映了一组数据的\textbf{中心趋势}（central tendency）。

在上面的问题中，若使两个射手各射 $N$ 枪，则他们打中的环数大约是：
\begin{align*}
\text{甲}&:8\times 0.3N+9\times 0.1N+10\times 0.6N=9.3N,\\
\text{乙}&:8\times 0.2N+9\times 0.5N+10\times 0.3N=9.1N.
\end{align*}
平均起来甲每枪射中 $9.3$ 环，乙射中 $9.1$ 环，所以甲射手的技术要好些。

因此，这里的计算公式为
\[
m=8\times p_8+9\times p_9+10\times p_{10},
\]
这是对击中环数的加权平均，其权正好取为相应的概率。

\subsection*{二、离散型场合}

受上面的问题启发，对一般离散型随机变量，我们可引进如下定义。

\begin{definition}[离散型随机变量的数学期望]\label{definition:ch04-离散型期望}
设 $\xi$ 为一离散型随机变量，它取值 $x_1,x_2,x_3,\cdots$，对应的概率为 $p_1,p_2,p_3,\cdots$。如果级数
\begin{align}
\sum_{i=1}^{\infty} x_i p_i\label{align:ch04-离散期望级数}
\end{align}
\textbf{绝对收敛}，则把它称为 $\xi$ 的\textbf{数学期望}（mathematical expectation），简称\textbf{期望}、\textbf{期望值}或\textbf{均值}（mean），记作 $E\xi$。当 $\sum_{i=1}^{\infty}|x_i|p_i$ 发散时，则说 $\xi$ 的数学期望不存在。
\end{definition}

定义中对级数要求绝对收敛是为了数学处理的方便。从直观上来讲，它也是合理的：因为诸 $x_i$ 的顺序对随机变量并不是本质的，因而在数学期望的定义中就应允许任意改变 $x_i$ 的次序而不影响其收敛性及其和值，这在数学上就相当于要求级数~\eqref{align:ch04-离散期望级数}绝对收敛。

显然数学期望由概率分布唯一确定，以后我们也称它为某概率分布的数学期望。下面来计算一些重要的离散型分布的期望值。

\begin{example}[Bernoulli 分布]\label{example:ch04-期望-Bernoulli}
事件 $A$ 发生的概率为 $p$，若以 $\mathbf{1}_A$ 记其示性函数，即 $A$ 发生时取值 $1$，否则取值 $0$，则
\begin{align}
E\mathbf{1}_A=1\times p+0\times(1-p)=p=P(A).\label{align:ch04-示性函数期望}
\end{align}
因此概率 $P(A)$ 是随机变量 $\mathbf{1}_A$ 的数学期望。从这个角度看，概率是数学期望的特例。
\end{example}

\begin{example}[二项分布的期望]\label{example:ch04-期望-二项}
二项分布 $p_k=\binom{n}{k}p^k q^{n-k}$，$k=0,1,2,\cdots,n$，则
\begin{align*}
\sum_{k=0}^n kp_k
&=\sum_{k=1}^n k\binom{n}{k}p^k q^{n-k}\\
&=np\sum_{k=1}^n \binom{n-1}{k-1}p^{k-1}q^{n-k}\\
&=np(p+q)^{n-1}=np.
\end{align*}
掷一枚均匀硬币 $100$ 次，您能期望得到多少次正面？——答案就是 $np$。
\end{example}

\begin{example}[Poisson 分布的期望]\label{example:ch04-期望-Poisson}
Poisson 分布 $p_k=\dfrac{\lambda^k}{k!}\mathrm{e}^{-\lambda}$，$k=0,1,2,\cdots$，则
\begin{align}
\sum_{k=0}^{\infty} kp_k
&=\sum_{k=1}^{\infty} k\cdot\frac{\lambda^k}{k!}\mathrm{e}^{-\lambda}
=\lambda\mathrm{e}^{-\lambda}\sum_{k=1}^{\infty}\frac{\lambda^{k-1}}{(k-1)!}\nonumber\\
&=\lambda\mathrm{e}^{-\lambda}\cdot\mathrm{e}^{\lambda}=\lambda.\label{align:ch04-Poisson期望}
\end{align}
由此看出，Poisson 分布的参数 $\lambda$ 就是它的期望值。
\end{example}

\begin{example}[几何分布的期望]\label{example:ch04-期望-几何}
几何分布 $p_k=q^{k-1}p$，$k=1,2,\cdots$，则
\begin{align}
\sum_{k=1}^{\infty} kp_k
&=\sum_{k=1}^{\infty} kq^{k-1}p=p(1+2q+3q^2+\cdots)\nonumber\\
&=p(q+q^2+q^3+\cdots)'=p\left(\frac{q}{1-q}\right)'\nonumber\\
&=p\frac{1}{(1-q)^2}=\frac{1}{p}.\label{align:ch04-几何期望}
\end{align}
重复抛一粒骰子，平均抛多少次才能出现一次么点？——假如您心中的答案是 $6$，那么这就是 $\dfrac{1}{p}$。
\end{example}

\begin{example}[期望不存在的离散分布]\label{example:ch04-期望不存在}
随机变量 $\xi$ 取值 $x_k=(-1)^k\dfrac{2^k}{k}$，$k=1,2,\cdots$，对应的概率为 $p_k=\dfrac{1}{2^k}$。则由于 $p_k\geqslant 0$，$\sum_{k=1}^{\infty}p_k=1$，因此它是概率分布，而且
\[
\sum_{k=1}^{\infty} x_k p_k=\sum_{k=1}^{\infty}(-1)^k\frac{1}{k}=-\ln 2.
\]
但由于
\[
\sum_{k=1}^{\infty}|x_k|p_k=\sum_{k=1}^{\infty}\frac{1}{k}=\infty,
\]
因此按定义 $\xi$ 的数学期望不存在。
\end{example}

从上面例子看到，几种重要的离散型分布，其参数都可由数学期望算得，因此数学期望是一个重要的概念。

\subsection*{三、应用实例}

从上面讨论中可以看出，数学期望刻画了随机变量取值的某种平均，有明显的直观含意。在许多问题中，数学期望的概念甚至比概率和随机变量等概念更易被人理解和接受。在概率论发展的早期，人们常用的概念和工具是数学期望。下面的例子将有助读者对这个重要概念的深入理解。

\begin{example}[押宝]\label{example:ch04-押宝}
押宝是赌博的一种，它以种种形式在世界各地流行，吞夺大量财富，现举一例以揭穿其本质。

在我国南方流行一种称为``捉水鸡''的押宝，其规则如下：由庄家摸出一只棋子，放在密闭的盒中，这只棋子可以是红的或黑的将、士、象、车、马、炮之一。赌客们把钱押在一块写有上述十二个字（六个红字、六个黑字）的台面的某个字上。押定后，庄家揭开盒子露出那只棋子。凡押中者（字和颜色都对）以 $1$ 比 $10$ 得到赏金，不中者其押金归庄家。

为对这种押宝的实质有个了解，最好考察一个赌徒当他押上 $1$ 元赌注之后的期望所得。显然其分布列为
\begin{center}
\begin{tabular}{c|cc}
$\xi$ & $0$ & $11$\\
\hline
$P$ & $\dfrac{11}{12}$ & $\dfrac{1}{12}$
\end{tabular}
\end{center}
因此其数学期望为 $\dfrac{11}{12}$。

由于支出（$1$ 元）和期望收入 $\left(\dfrac{11}{12}\text{元}\right)$ 不等，因此这是不公平的赌博，它明显对庄家有利。数学期望的概念帮助我们认清了这个本质。

事实上，当一个赌徒走进一座赌场时，他面临的都是这种不公平的赌博，否则赌场的巨大开销和业主的高额利润从何而来呢？

普及概率知识有助于杜绝赌博现象。
\end{example}

\begin{example}[彩票]\label{example:ch04-彩票}
彩票的发行，数额巨大，其实质如何呢？请看一则实例：发行彩票 $100$ 万张，每张 $5$ 元。设头等奖 $5$ 个，奖金 $31.5$ 万元；二等奖 $95$ 个，奖金各 $5\,000$ 元；三等奖 $900$ 个，奖金各 $300$ 元；四等奖 $9\,000$ 个，奖金各 $20$ 元。

还是算算每张彩票的期望所得。这时分布列为
\begin{center}
\begin{tabular}{c|ccccc}
$\eta$ & $315\,000$ & $5\,000$ & $300$ & $20$ & $0$\\
\hline
$P$ & $\dfrac{5}{100\text{万}}$ & $\dfrac{95}{100\text{万}}$ & $\dfrac{900}{100\text{万}}$ & $\dfrac{9\,000}{100\text{万}}$ & $*$
\end{tabular}
\end{center}
其中一等奖的金额本来另行摇出，此地为简便计，用其均值；至于 $*$，无需细算。

花 $5$ 元买来的一张彩票，从摇奖中的期望所得为
\begin{align*}
E\eta&=315\,000\times\frac{5}{100\text{万}}+5\,000\times\frac{95}{100\text{万}}
+300\times\frac{900}{100\text{万}}+20\times\frac{9\,000}{100\text{万}}\\
&=2.5(\text{元}).
\end{align*}
即大约能收回一半。因此这实质上也是一种于购买者不利的非公平博弈，所以历来博彩并称。显然不能把购买彩票当作一种投资渠道。

在我国，彩票的发行严格由民政部门管理，只有当收益主要用于公益事业时才允许，如福利彩票与体育彩票。
\end{example}

\begin{example}[投资之决策]\label{example:ch04-投资决策}
投资总具有一定风险，因此在选择投资方向时，计算其期望收益常是可供考虑的决策方法之一。下面是一个大为简化的例子。

某人有 $10$ 万元现金，想投资于某项目，预估成功的机会为 $30\%$，可得利润 $8$ 万元，失败的机会为 $70\%$，将损失 $2$ 万元。若存入银行，同期间的利率为 $5\%$，问是否应作此项投资？

以 $\zeta$ 记投资利润，则其分布列为
\begin{center}
\begin{tabular}{c|cc}
$\zeta$ & $8$ & $-2$\\
\hline
$P$ & $0.3$ & $0.7$
\end{tabular}
\end{center}
因此
\[
E\zeta=8\times 0.3-2\times 0.7=1(\text{万元}),
\]
而存入银行的利息为 $10\times 5\%=0.5$（万元），因此从期望收益的角度看，似应选择投资，当然这时要看投资者是否愿意冒这样的风险。

投资总希望得到尽可能高的收益，又想尽量规避风险，问题是如何折衷其间。在引进风险的度量之后，我们将再次回到这个问题。

上述对赌博、彩票和投资的讨论集中于分析未来不确定性事件的金融后果，也彰显了数学期望概念的重要性。类似的还有保险。保险与赌博等不同，因为赌博创造一个本来没有的风险，而保险则管理一个本已存在且不可避免的风险。
\end{example}

\begin{example}[保险]\label{example:ch04-保险}
为规避某种未来的不确定性事件（例如空难），人们不是自己创造一个分担群，而是找一个中介即保险公司，付保险金。保险公司创造满足某些要求的群体，使其整个索赔费用为合理而可预测的（主要依据是下章要讲的大数定律），并在这个过程中得到利润。下面我们讨论如何对保险费进行精算。在保险学中，收取保险费的原则是：被保险人交的``纯保险费''与他们所能得到的赔偿金的期望值相等。

因此，若出事的概率为 $p$，有 $N$ 个人参加保险，则每人交的纯保险费 $a$ 与出事赔偿金 $b$，应有下面关系：
\[
Na=\sum_{k=0}^N \binom{N}{k}p^k(1-p)^{N-k}\cdot kb,
\]
即 $a=pb$，正如所预期的。

保险的种类繁多，但上述原则不变，因此保险费的计算离不开概率论。
\end{example}

\begin{example}[一种验血新技术]\label{example:ch04-验血}
在一个人数很多的单位中普查某种疾病，$N$ 个人去验血。对这些人的血的化验可以用两种办法进行。(1) 每个人的血分别化验，这时需要化验 $N$ 次；(2) 把 $k$ 个人的血混在一起进行化验。如果结果是阴性的，那么对这 $k$ 个人只作一次检验就够了；如果结果是阳性的，那么必须对这 $k$ 个人再逐个分别化验，这时对这 $k$ 个人共需作 $k+1$ 次化验。假定对所有的人来说，化验是阳性反应的概率都是 $p$，而且这些人的反应是独立的。我们来说明在 $p$ 相当小的场合，采用办法 (2) 能减少化验的次数。

若记 $q=1-p$，则 $k$ 个人的混血呈阳性反应的概率为 $1-q^k$。用 (2) 的方法验血时，每个人的血需要化验的次数 $\xi$ 是随机变量，其分布列为
\begin{center}
\begin{tabular}{c|cc}
$\xi$ & $\dfrac{1}{k}$ & $1+\dfrac{1}{k}$\\
\hline
$P$ & $q^k$ & $1-q^k$
\end{tabular}
\end{center}
因此
\begin{align*}
E\xi=\frac{1}{k}\cdot q^k+\left(1+\frac{1}{k}\right)(1-q^k)=1-q^k+\frac{1}{k}.
\end{align*}
$N$ 个人需要的化验次数的期望值为 $N\left(1-q^k+\dfrac{1}{k}\right)$，当 $q^k-\dfrac{1}{k}>0$ 时，就能减少验血次数。例如当 $p=0.1$ 时，取 $k=4$，则 $q^k-\dfrac{1}{k}=0.4$。用 (2) 法平均能减少 $40\%$ 的工作量。显然 $p$ 愈小，用这种方法愈有利。当 $p$ 已知时，还可以选定整数 $k_0$，使 $E\xi$ 达到最小，把 $k_0$ 个人分为一组就最能节省化验次数。
\end{example}

\subsection*{四、连续型场合}

下面我们转入考虑连续型随机变量的数学期望。设随机变量 $\xi$ 有密度函数 $p(x)$，取很密的分点 $x_0<x_1<x_2<\cdots<x_n$，则 $\xi$ 落在 $[x_i,x_{i+1})$ 中的概率近似地等于 $p(x_i)(x_{i+1}-x_i)$，因此 $\xi$ 与以概率 $p(x_i)(x_{i+1}-x_i)$ 取值 $x_i$ 的离散型随机变量近似，而这离散型随机变量的数学期望为
\[
\sum_i x_i p(x_i)(x_{i+1}-x_i).
\]
上式是积分 $\int_{-\infty}^{\infty} xp(x)\,\mathrm{d}x$ 的渐近和式，这个直观的考虑启发我们引进如下定义。

\begin{definition}[连续型随机变量的数学期望]\label{definition:ch04-连续型期望}
设 $\xi$ 为具有密度函数 $p(x)$ 的连续型随机变量，当积分 $\int_{-\infty}^{\infty} xp(x)\,\mathrm{d}x$ 绝对收敛时，我们称它为 $\xi$ 的\textbf{数学期望}（或均值），记作 $E\xi$，即
\begin{align}
E\xi=\int_{-\infty}^{\infty} xp(x)\,\mathrm{d}x.\label{align:ch04-连续期望}
\end{align}
\end{definition}

显然这里定义的数学期望也只与分布有关。下面计算一些重要的连续型分布的数学期望。

\begin{example}[正态分布的期望]\label{example:ch04-期望-正态}
正态分布 $N(\mu,\sigma^2)$，则
\begin{align}
\int_{-\infty}^{\infty} xp(x)\,\mathrm{d}x
&=\int_{-\infty}^{\infty} x\frac{1}{\sqrt{2\pi}\sigma}\mathrm{e}^{-(x-\mu)^2/(2\sigma^2)}\,\mathrm{d}x\nonumber\\
&=\frac{1}{\sqrt{2\pi}}\int_{-\infty}^{\infty}(\sigma z+\mu)\mathrm{e}^{-z^2/2}\,\mathrm{d}z\nonumber\\
&=\frac{\mu}{\sqrt{2\pi}}\int_{-\infty}^{\infty}\mathrm{e}^{-z^2/2}\,\mathrm{d}z=\mu.\label{align:ch04-正态期望}
\end{align}
可见 $N(\mu,\sigma^2)$ 中的参数 $\mu$ 正是它的数学期望。
\end{example}

\begin{example}[指数分布的期望]\label{example:ch04-期望-指数}
指数分布 $p(x)=\lambda\mathrm{e}^{-\lambda x}$，$x\geqslant 0$，则
\begin{align}
\int_0^{\infty} x\lambda\mathrm{e}^{-\lambda x}\,\mathrm{d}x
=-\int_0^{\infty} x\,\mathrm{d}\mathrm{e}^{-\lambda x}
=\int_0^{\infty}\mathrm{e}^{-\lambda x}\,\mathrm{d}x=\frac{1}{\lambda}.\label{align:ch04-指数期望}
\end{align}
\end{example}

\begin{example}[Cauchy 分布的期望不存在]\label{example:ch04-期望-Cauchy}
Cauchy 分布 $p(x)=\dfrac{1}{\pi}\cdot\dfrac{1}{1+x^2}$。由于
\[
\int_{-\infty}^{\infty}|x|\cdot\frac{1}{\pi(1+x^2)}\,\mathrm{d}x=\infty,
\]
因此 Cauchy 分布的数学期望不存在。
\end{example}

\subsection*{五、一般场合}

我们已经对离散型随机变量及连续型随机变量分别定义了数学期望，现在自然希望找到一种能适合一切随机变量的数学期望定义，并把上述两种情况作为特例。为了做到这点，需要利用 Stieltjes 积分\footnote{如果读者不熟悉 Stieltjes 积分而又希望对它有所了解，可参看复旦大学数学系《实变数函数论与泛函分析概要》（第二版），上海科技出版社，1963，第四章。不过，为读懂本书其余部分，只要承认后面所述的少数事实就可以了。}。

若随机变量 $\xi$ 的分布函数为 $F(x)$，类似于连续型随机变量的场合，作很密的分割 $x_0<x_1<x_2<\cdots<x_n$，则 $\xi$ 落在 $[x_i,x_{i+1})$ 中的概率等于 $F(x_{i+1})-F(x_i)$，因此 $\xi$ 与以概率 $F(x_{i+1})-F(x_i)$ 取值 $x_i$ 的离散型随机变量近似，而后者的数学期望为
\[
\sum_i x_i[F(x_{i+1})-F(x_i)].
\]
注意到上式是 Stieltjes 积分 $\int_{-\infty}^{\infty} x\,\mathrm{d}F(x)$ 的渐近和式，这启发我们引进如下定义。

\begin{definition}[一般随机变量的数学期望]\label{definition:ch04-Stieltjes期望}
若 $\xi$ 的分布函数为 $F(x)$，则定义
\begin{align}
E\xi=\int_{-\infty}^{\infty} x\,\mathrm{d}F(x)\label{align:ch04-Stieltjes期望}
\end{align}
为 $\xi$ 的\textbf{数学期望}（或均值）。这里我们还是要求上述积分绝对收敛，否则数学期望不存在。
\end{definition}

我们仅列举 Stieltjes 积分
\begin{align}
I=\int_{-\infty}^{\infty} g(x)\,\mathrm{d}F(x)\label{align:ch04-Stieltjes一般}
\end{align}
的如下性质：

(i) 当 $F(x)$ 为跳跃函数，在 $x_i$（$i=1,2,\cdots$）具有跃度 $p_i$ 时，上面积分化为无穷级数
\[
I=\sum_i g(x_i)p_i;
\]

(ii) 当 $F(x)$ 存在导数 $F'(x)=p(x)$ 时，积分~\eqref{align:ch04-Stieltjes一般}化为普通积分
\[
I=\int_{-\infty}^{\infty} g(x)p(x)\,\mathrm{d}x;
\]

(iii) 线性性质
\[
\int_{-\infty}^{\infty}[ag_1(x)+bg_2(x)]\,\mathrm{d}F(x)
=a\int_{-\infty}^{\infty}g_1(x)\,\mathrm{d}F(x)+b\int_{-\infty}^{\infty}g_2(x)\,\mathrm{d}F(x);
\]

(iv)
\[
\int_{-\infty}^{\infty}g(x)\,\mathrm{d}[aF_1(x)+bF_2(x)]
=a\int_{-\infty}^{\infty}g(x)\,\mathrm{d}F_1(x)+b\int_{-\infty}^{\infty}g(x)\,\mathrm{d}F_2(x);
\]

(v)
\[
\int_a^b g(x)\,\mathrm{d}F(x)=\int_a^c g(x)\,\mathrm{d}F(x)+\int_c^b g(x)\,\mathrm{d}F(x),\qquad(a\leqslant c\leqslant b);
\]

(vi) 若 $g(x)\geqslant 0$，$F(x)$ 单调不减，$b>a$，则
\[
\int_a^b g(x)\,\mathrm{d}F(x)\geqslant 0.
\]

从头两个性质，我们知道定义式~\eqref{align:ch04-Stieltjes期望}的确能包含~\eqref{align:ch04-离散期望级数}及~\eqref{align:ch04-连续期望}作为特例。

\subsection*{六、随机变量函数的数学期望}

下面讨论随机变量的函数 $\eta=g(\xi)$ 的数学期望的定义。这里 $\xi$ 是分布函数为 $F_\xi(x)$ 的随机变量，而 $g(x)$ 是一元 Borel 函数，从上章 $\S 3$ 的讨论已知 $\eta$ 是随机变量。通过类似于引进~\eqref{align:ch04-Stieltjes期望}的推理，似应定义 $g(\xi)$ 的数学期望为
\[
\sum_i g(x_i)[F_\xi(x_{i+1})-F_\xi(x_i)]
\]
的极限，即
\begin{align}
Eg(\xi)=\int_{-\infty}^{\infty} g(x)\,\mathrm{d}F_\xi(x).\label{align:ch04-随机变量函数期望}
\end{align}
这里当然还是要求这个积分绝对收敛。

但是，另一方面，因为 $\eta$ 是随机变量，也有分布函数，记之为 $F_\eta(x)$，则按一般随机变量数学期望的定义式~\eqref{align:ch04-Stieltjes期望}又应有
\begin{align}
E\eta=\int_{-\infty}^{\infty} y\,\mathrm{d}F_\eta(y).\label{align:ch04-eta期望}
\end{align}
因此，这两个积分应该相等。事实上，这两个积分的确相等，我们把这个事实写成定理的形式。

\begin{theorem}[佚名统计学家公式]\label{theorem:ch04-佚名统计学家公式}
若 $g(x)$ 是一元 Borel 函数，而 $\eta=g(\xi)$，则
\begin{align}
\int_{-\infty}^{\infty} y\,\mathrm{d}F_\eta(y)=\int_{-\infty}^{\infty} g(x)\,\mathrm{d}F_\xi(x).\label{align:ch04-两个积分相等}
\end{align}
即这两个积分中，若有一个存在，则另一个也存在，而且两者相等。
\end{theorem}

这个定理的证明要用到测度论，超出了本课程范围。我们只能列出结论，并指出它的重要性：一方面，它消除了随机变量的数学期望定义中所出现的表面矛盾；另一方面在计算随机变量函数的数学期望时也可以带来很大的方便，我们无须先计算 $\eta$ 的分布函数 $F_\eta(y)$（回忆上章 $\S 3$ 的繁复计算）再求其数学期望，而可以直接从 $\xi$ 的分布函数 $F_\xi(x)$ 出发利用~\eqref{align:ch04-随机变量函数期望}来计算。

历史上，先是统计学家从直观出发广泛使用~\eqref{align:ch04-随机变量函数期望}式，尔后才由概率论专家从数学上给予严格证明，首创者已难追寻，姑且名为\textbf{佚名统计学家公式}。

佚名统计学家公式~\eqref{align:ch04-随机变量函数期望}有十分明显的直观解释，在离散型场合，它化为
\begin{align}
Eg(\xi)=\sum_{i=1}^{\infty} g(x_i)p(x_i),\label{align:ch04-离散函数期望}
\end{align}
这从 $g(\xi)$ 的分布列~(3.3.5) 立刻就能得到。

我们以一个浅显的例子来说明它的含意。

某零售商店出售某种小商品，每销售一件可赚 $1.5$ 元，关心的是它的利润 $\eta$。由于利润与销售量 $\xi$ 成正比，因此要知道销售量。比较合理的做法是假定 $\xi$ 为随机变量，服从一定分布。为简单起见，假定每天卖出 $0,1,2,3$ 件的概率分别为 $0.4,0.3,0.2,0.1$。这时希望计算平均利润 $E\eta$。

第一种计算方法是把 $\eta$ 看作随机变量，求出其概率分布，再按一般数学期望公式计算。这时 $\eta$ 的分布列为
\begin{center}
\begin{tabular}{c|cccc}
$\eta$ & $0$ & $1.5$ & $3$ & $4.5$\\
\hline
$P$ & $0.4$ & $0.3$ & $0.2$ & $0.1$
\end{tabular}
\end{center}
因此
\begin{align*}
E\eta&=0\times 0.4+1.5\times 0.3+3\times 0.2+4.5\times 0.1\\
&=1.5(\text{元}).
\end{align*}
这种算法就是按公式~\eqref{align:ch04-eta期望}的思路进行的。

另一种算法把 $\eta$ 看作 $\xi$ 的函数 $\eta=1.5\xi$，由于 $\xi$ 的分布列早就知道为
\begin{center}
\begin{tabular}{c|cccc}
$\xi$ & $0$ & $1$ & $2$ & $3$\\
\hline
$P$ & $0.4$ & $0.3$ & $0.2$ & $0.1$
\end{tabular}
\end{center}
因此按~\eqref{align:ch04-离散函数期望}式（也即按~\eqref{align:ch04-随机变量函数期望}式）
\begin{align*}
E\eta=E(1.5\xi)&=1.5\times 0\times 0.4+1.5\times 1\times 0.3+1.5\times 2\times 0.2+1.5\times 3\times 0.1\\
&=1.5(\text{元}).
\end{align*}
正如预期的，两个答案相同，这也为\cref{theorem:ch04-佚名统计学家公式}提供了例证。

下面的例子更接近于实际，不过其方法实质是一样的。

\begin{example}[报童问题]\label{example:ch04-报童问题}
设某报童每日的潜在卖报数 $\zeta$ 服从参数为 $\lambda$ 的 Poisson 分布。如果每卖出一份报可得报酬 $a$，卖不掉而退回则每份赔偿 $b$。若某日该报童批进 $n$ 份报，试求其期望所得。进一步，还要求最佳的批进份数 $n$。
\end{example}

\begin{proof}[解]
若记真正卖报数为 $\xi$，则 $\xi$ 与 $\zeta$ 的关系如下：
\[
\xi=\begin{cases}
\zeta,&\zeta<n,\\
n,&\zeta\geqslant n.
\end{cases}
\]
这里的 $\xi$ 服从截尾 Poisson 分布，即
\[
P\{\xi=k\}=
\begin{cases}
\dfrac{\lambda^k}{k!}\mathrm{e}^{-\lambda},&k<n,\\[6pt]
\displaystyle\sum_{i=n}^{\infty}\frac{\lambda^i}{i!}\mathrm{e}^{-\lambda},&k=n.
\end{cases}
\]

记所得为 $\eta$，则 $\eta$ 与 $\xi$ 的关系如下：
\[
\eta=g(\xi)=
\begin{cases}
an,&\xi=n,\\
a\xi-b(n-\xi),&\xi<n.
\end{cases}
\]
因此由~\eqref{align:ch04-离散函数期望}，期望所得为
\begin{align*}
M(n)=Eg(\xi)
&=\sum_{k=0}^{n-1}\frac{\lambda^k}{k!}\mathrm{e}^{-\lambda}\bigl[ka-(n-k)b\bigr]
+\left(\sum_{k=n}^{\infty}\frac{\lambda^k}{k!}\mathrm{e}^{-\lambda}\right)na\\
&=(a+b)\lambda\sum_{k=0}^{n-2}\frac{\lambda^k}{k!}\mathrm{e}^{-\lambda}
-n(a+b)\sum_{k=0}^{n-1}\frac{\lambda^k}{k!}\mathrm{e}^{-\lambda}+na.
\end{align*}

这个问题的最终解决是当 $a,b,\lambda$ 给定后，求 $n$ 使 $M(n)$ 达到极大，这是一个典型的最优化问题，也是报童问题的正确解。

由~(3.1.44) 推知，计算 Poisson 分布的部分和可用下列公式：
\begin{align}
\sum_{k=r}^{\infty}\frac{\lambda^k}{k!}\mathrm{e}^{-\lambda}
=\int_0^1\frac{\lambda^r}{\Gamma(r)}x^{r-1}\mathrm{e}^{-\lambda x}\,\mathrm{d}x
=\frac{1}{\Gamma(r)}\int_0^{\lambda}y^{r-1}\mathrm{e}^{-y}\,\mathrm{d}y.\label{align:ch04-Poisson部分和}
\end{align}
右端的数值可在数学表中查到。公式~\eqref{align:ch04-Poisson部分和} 是 Erlang 分布的一个应用，也可用分析方法直接证明。
\end{proof}

佚名统计学家公式更有用的场合是当 $\xi$ 具有密度函数 $p(x)$ 时，这时它化为
\begin{align}
E\eta=Eg(\xi)=\int_{-\infty}^{\infty} g(x)p(x)\,\mathrm{d}x.\label{align:ch04-连续函数期望}
\end{align}
这个公式经常用到。下面是一个例子，这个例子也顺带说明可以利用随机变量的期望值来作出某种最优决策。

\begin{example}[最优存货量]\label{example:ch04-最优存货量}
市场上对某种商品的需求量是随机变量 $\xi$（单位：吨），它服从 $[2\,000,4\,000]$ 上均匀分布，设每售出这种商品 $1$ 吨，可挣得 $3$ 万元，但假如销售不出而屯积于仓库，则每吨需浪费保养费 $1$ 万元，问题是要确定应组织多少货源，才能使收益最大？
\end{example}

\begin{proof}[解]
若以 $y$ 记预备的此种商品量（显然可以只考虑 $2\,000\leqslant y\leqslant 4\,000$ 的情况），则收益（单位：万元）
\[
\eta=H(\xi)=\begin{cases}
3y,&\text{当 }\xi\geqslant y\text{ 时},\\
3\xi-(y-\xi),&\text{当 }\xi<y\text{ 时}.
\end{cases}
\]
为了求得 $E\eta$，利用公式~\eqref{align:ch04-连续函数期望}，
\begin{align*}
E\eta=\int_{-\infty}^{\infty}H(x)p(x)\,\mathrm{d}x
&=\frac{1}{2\,000}\int_{2\,000}^{4\,000}H(x)\,\mathrm{d}x\\
&=\frac{1}{2\,000}\int_{2\,000}^{y}(4x-y)\,\mathrm{d}x+\frac{1}{2\,000}\int_y^{4\,000}3y\,\mathrm{d}x\\
&=\frac{1}{1\,000}\bigl(-y^2+7\,000y-4\times 10^6\bigr).
\end{align*}
此式当 $y=3\,500$ 时达到最大，因此组织 $3\,500$ 吨此种商品是最好的决策。
\end{proof}

\subsection*{七、多维场合}

可以把佚名统计学家公式推广到随机向量的场合，若 $(\xi_1,\cdots,\xi_n)$ 的分布函数为 $F(x_1,\cdots,x_n)$，而 $g(x_1,\cdots,x_n)$ 为 $n$ 元 Borel 函数，则
\[
Eg(\xi_1,\cdots,\xi_n)=\int_{-\infty}^{\infty}\cdots\int_{-\infty}^{\infty}g(x_1,\cdots,x_n)\,\mathrm{d}F(x_1,\cdots,x_n).
\]
特别地，
\[
E\xi_1=\int_{-\infty}^{\infty}\cdots\int_{-\infty}^{\infty}x_1\,\mathrm{d}F(x_1,\cdots,x_n)
=\int_{-\infty}^{\infty}x_1\,\mathrm{d}F_1(x_1),
\]
其中 $F_1(x_1)$ 是 $\xi_1$ 的分布函数。一般地，引进如下定义。

\begin{definition}[随机向量的数学期望]\label{definition:ch04-随机向量期望}
随机向量 $(\xi_1,\xi_2,\cdots,\xi_n)$ 的\textbf{数学期望}为 $(E\xi_1,E\xi_2,\cdots,E\xi_n)$，其中
\[
E\xi_i=\int_{-\infty}^{\infty}\cdots\int_{-\infty}^{\infty}x_i\,\mathrm{d}F(x_1,\cdots,x_n)
=\int_{-\infty}^{\infty}x_i\,\mathrm{d}F_i(x_i),
\]
这里 $F_i(x_i)$ 是 $\xi_i$ 的分布函数（$i=1,2,\cdots,n$）。
\end{definition}

因此二元正态分布~(3.2.22) 中的参数 $(\mu_1,\mu_2)$ 正是相应的随机向量的数学期望。

一维或多维的佚名统计学家公式在以下的理论推导与数值计算中都起重要作用，它是概率论中常用公式之一。

\subsection*{八、数学期望的基本性质}

\noindent\textbf{性质 1（期望的有界性）}\par
若 $a\leqslant\xi\leqslant b$，则 $a\leqslant E\xi\leqslant b$。特别地 $Ec=c$，这里 $a,b,c$ 是常数。


\noindent\textbf{性质 2（期望的线性性质）}\par
对任意常数 $c_i$，$i=1,\cdots,n$ 及 $b$，有
\[
E\left(\sum_{i=1}^n c_i\xi_i+b\right)=\sum_{i=1}^n c_i E\xi_i+b.
\]


利用佚名统计学家公式，这两个性质的证明是明显的。期望的线性性质 对数学期望的计算很有用处。

\begin{example}[超几何分布的期望]\label{example:ch04-超几何期望}
求超几何分布
\[
p_m=\frac{\binom{M}{m}\binom{N-M}{n-m}}{\binom{N}{n}},\qquad m=0,1,\cdots,n
\]
的数学期望。
\end{example}

当然可以用 $\sum_{k=0}^n kp_k$ 直接求出，但也可用下面方法来计算。设想一个相应的不放回抽样，令
\[
\xi_i=\begin{cases}
1,&\text{第 }i\text{ 次抽得次品},\\
0,&\text{第 }i\text{ 次抽得好品}.
\end{cases}
\]
则 $P\{\xi_i=1\}=\dfrac{M}{N}$，因此 $E\xi_i=\dfrac{M}{N}$，而 $\xi=\xi_1+\cdots+\xi_n$ 表示 $n$ 次不放回抽样中抽出的次品数，它服从上述超几何分布，利用期望的线性性质得到
\[
E\xi=E\xi_1+\cdots+E\xi_n=\frac{nM}{N}.
\]

\end{document}